% Generated draft: requires mathematical and source review.
\documentclass[11pt]{article}
\usepackage[margin=25mm]{geometry}
\usepackage{fontspec}
\setmainfont{Latin Modern Roman}
\usepackage{amsmath,amssymb,mathtools,unicode-math}
\setmathfont{Latin Modern Math}
\usepackage{enumitem,needspace,xurl,hyperref}
\setlength{\parindent}{0pt}
\setlength{\parskip}{5pt}
\setlength{\emergencystretch}{3em}
\newcommand{\N}{\mathbb{N}}
\newcommand{\Z}{\mathbb{Z}}
\newcommand{\Q}{\mathbb{Q}}
\newcommand{\R}{\mathbb{R}}
\newcommand{\C}{\mathbb{C}}
\newcommand{\F}{\mathbb{F}}
\newcommand{\A}{\mathbb{A}}
\newcommand{\PP}{\mathbb{P}}
\newcommand{\E}{\mathbb{E}}
\newcommand{\eps}{\varepsilon}
\newcommand{\dd}{\,\mathrm{d}}
\begin{document}
\section*{UnsolvedMath ID: 11000230}
\textbf{Source problem code:} catalogue entry
\subsection{Definitions and mathematical statement}
The problem asks whether the converse of the preceding theorem remains true after dropping its assumption that \(G\) is free. That theorem is absent, so its other hypotheses and the exact implication being reversed cannot be recovered uniquely. The surrounding programme concerns word hyperbolicity of surface-group extensions. In the customary notation, such an extension has the form
\[1\longrightarrow\pi_1(S)\longrightarrow\Gamma_G\longrightarrow G\longrightarrow1,\]
with the action/extension data specified by the source; this excerpt does not give their construction or a particular property of \(G\) to substitute for the missing theorem.

The stated goal is to extend the Bestvina--Feighn combination principle from actions on trees to actions on Gromov-hyperbolic cell complexes, with sufficient additional structure to decide word hyperbolicity of \(\Gamma_G\) in the intended setting. A finitely generated group is word hyperbolic when a Cayley graph has uniformly thin geodesic triangles. An action on a hyperbolic complex by itself is not the asserted criterion; the missing hypotheses must specify how stabilizers, attachment maps and extension geometry are controlled.

The source specifically points to the action of \(\Gamma_G\) on a Rips complex of \(G\). Given a word metric on \(G\) and scale \(R\), that complex has vertices \(G\) and a simplex on each finite subset of diameter at most \(R\), with the translation action of \(G\); the quotient map \(\Gamma_G\to G\) induces the indicated action. The kernel acts trivially in this model, so no proper action or immediate geometric-action criterion is inferred. The metric, scale, hyperbolicity and finite-generation assumptions on \(G\) are not stated in the extract.

The source says there were no available examples on which to test its Problem 2.2, and announces related proposals in Section 4. This is a historical limitation of that investigation, not a current nonexistence assertion. Section 4 and the preceding theorem are not supplied. In particular, a familiar convex-cocompactness or other surface-subgroup theorem is not silently chosen as the missing converse.
\subsection{Short English statement}
Extend the omitted surface-extension converse beyond free quotient groups, using a combination theorem for actions on hyperbolic cell complexes rather than only trees. The exact theorem and its remaining hypotheses are missing.
\subsection{Sources}
\begin{itemize}
\item Local UnsolvedMath registry, permanent ID \(11000230\), exact source code \texttt{catalogue entry}.
\item Farb, \emph{Problems on Mapping Class Groups and Related Topics}, Problem 2.2, PDF page 269; running chapter title \emph{Problems in the geometry of surface group extensions}: \url{https://www.math.uchicago.edu/~farb/papers/mcgbook.pdf}. Registry access 2026-07-29; pdf-font-structure extraction.
\item Bestvina--Feighn's combination theorem is named without a full citation; the preceding theorem and Section 4 proposals are absent.
\end{itemize}
\end{document}
